\subsection*{B.1 Entrevista N.°~1 --- Delfina Fernández}
\addcontentsline{toc}{subsection}{B.1 Entrevista N.° 1 --- Delfina Fernández}

\noindent\textbf{Entrevistada:} Delfina Fernández\\
\noindent\textbf{Perfil:} Dueña de emprendimiento de indumentaria bordada a mano (remeras y musculosas). Vende por Instagram y Tienda Nube. También trabaja en armado de pedidos de tienda online (otra empresa).\\
\noindent\textbf{Entrevistadores:} Tomás Gonçalves (moderador) y Micaela Palomino (observadora)\\
\noindent\textbf{Fecha:} Domingo 31 de mayo de 2026\\
\noindent\textbf{Modalidad:} Videollamada (semiestructurada)\\
\noindent\textbf{Plataforma del negocio:} Tienda Nube (primer nivel pago)

\begin{framed}
\small\noindent\textit{Nota metodológica: Esta transcripción fue limpiada a partir de la grabación original. Se eliminó charla trivial previa y posterior a la entrevista, se corrigió la sintaxis para mejorar la fluidez de lectura y se reasignaron las atribuciones de orador en casos donde la herramienta de transcripción automática las asignó erróneamente (verificado por contexto). No se agregó, inventó ni dedujo información que los participantes no hayan dicho explícitamente. La entrevistada cuenta con un colaborador (Facundo) que interviene brevemente mostrando una aplicación personalizada.}
\end{framed}

\subsubsection*{Bloque 0 · Introducción y rapport}

\noindent\textbf{Micaela Palomino:} Nuestro proyecto, nuestra tesis, trata de hacer un agente y conectarlo a Tienda Nube. Me acordé de que vos tenés un emprendimiento en Tienda Nube, así que nos viniste a escuchar.

\noindent\textbf{Delfina Fernández:} Sí, yo te iba a decir, tengo un poco de miedo porque en realidad no sé qué tanto uso Tienda Nube. Pero igual trabajo con Tienda Nube; mi trabajo es con Tienda Nube. Así que bueno.

\noindent\textbf{Micaela Palomino:} La entrevista la va a llevar Tommy, así que todo tuyo, Tommy.

\noindent\textbf{Tomás Gonçalves:} Arranquemos de base, porque yo no conozco. Delfina, ¿de qué es tu negocio? ¿Qué vendés?

\noindent\textbf{Delfina Fernández:} Yo tengo un emprendimiento ---en realidad, porque es chiquito todavía--- en el que hago remeras y musculosas bordadas a mano. Por ahora es eso. Tengo mi página de Instagram y también mi página de Tienda Nube para vender por ahí. La verdad es que vendo por los dos lugares, pero Tienda Nube sirve mucho.

\subsubsection*{Bloque 1 · Contexto del negocio y herramientas actuales}

\noindent\textbf{Tomás Gonçalves:} Me acuerdo que dijiste que empezaste en Tienda Nube, pero como que no sabés usarlo tanto en la parte de datos.

\noindent\textbf{Delfina Fernández:} A ver, no es que tengo un máster en Tienda Nube. Con el emprendimiento no lo uso tanto. Las ventas, quizás por ahora, las estoy haciendo más en el boca a boca. Pero en mi trabajo, trabajo con todo lo que es armado de pedidos de tienda online. Entonces, ahí sí uso Tienda Nube, todo el tema de stock y demás.

\noindent\textbf{Tomás Gonçalves:} De la data que te ofrece Tienda Nube, ¿le das mucha pelota?

\noindent\textbf{Delfina Fernández:} Me imagino que tiene que tener estadísticas, tema de ventas. Eso por ahora no tanto, porque recién ahora empecé a pagar el primer nivel, que es el que te da más datos, más opciones para métodos de pago y métodos de envío. Todavía no lo uso tanto, pero porque tampoco lo tengo linkeado con Instagram, que creo que se puede hacer. Así que veo, por ejemplo, cuántas visitas tuve en el día, y esas cosas.

\noindent\textbf{Micaela Palomino:} Vos entrás en la Tienda Nube del emprendimiento. Básicamente lo que hacés es cargar el stock que tenés. ¿Y después cómo es todo el flujo? ¿Te llegan las ventas? ¿Se gestionan por Tienda Nube? ¿Vos tenés que entrar a ver qué ventas se realizaron, o te llega un mail?

\noindent\textbf{Delfina Fernández:} Me llegan las dos cosas. Cuando llega una venta me llegan tres notificaciones más un mail o dos, si el pago ya está recibido. Yo el stock ahí no lo tengo cargado todo, porque como son productos que tardo bastante en hacer, me da miedo cargar cinco remeras de un modelo y que me compren las cinco en la misma semana cuando ya tengo diez ventas y no las puedo entregar a tiempo. Entonces eso lo voy gestionando yo: pongo uno o dos de cada uno, y cuando veo que alguien me compra y ya no tengo stock de eso, lo doy de baja. Yo me manejo con las musculosas: quizás tengo la blanca en distintos modelos, pero en realidad tengo cinco modelos de bordado y dos remeras blancas nada más. Entonces, cuando me compran una, ya tengo que estar atenta a la próxima compra para darla de baja, porque no tengo otra.

\noindent\textbf{Micaela Palomino:} Y cuando te hacen una compra, ¿vos tenés que entrar a Tienda Nube a fijarte que te quedaste sin esto y volver a recargarlo? ¿Es todo manual, o te dice que te quedaste sin stock?

\noindent\textbf{Delfina Fernández:} Me lo dice. Las tres notificaciones que me llegan son: tenés una nueva compra (que más o menos te dice qué te compraron y quién), después me llega otra que dice que me quedé sin stock de ese producto exacto, y si ya pagaron, el pago recibido. Además me llega mail de la compra y mail del pago recibido.

\noindent\textbf{Micaela Palomino:} ¿Y te avisa antes de quedarte sin stock? Cuando te queda una prenda, ¿no te avisa?

\noindent\textbf{Delfina Fernández:} No, no, no. Eso no.

\noindent\textbf{Tomás Gonçalves:} O sea, es todo dinámico: te llega la alerta de que te quedaste sin stock y decís, bueno, me toca reponer.

\noindent\textbf{Delfina Fernández:} Claro, literalmente.

\noindent\textbf{Tomás Gonçalves:} ¿Y entre que te quedaste sin stock y reponés, cuánto tiempo pasa?

\noindent\textbf{Delfina Fernández:} Depende de cuándo hagan la compra. Si lo hacen durante el día, lo hago rápido: me fijo si sigo teniendo stock de eso y agrego otra más. Pero, por ejemplo, el otro día me desperté y había una compra de muy temprano; era sábado, como a las ocho de la mañana. Yo hasta las once no abrí los ojos, la vi, y ahí recién repuse el stock. Pero si no, es rápido.

\noindent\textit{--- Métricas y estadísticas de Tienda Nube ---}

\noindent\textbf{Micaela Palomino:} Me quedé con lo que decías de las visitas, que era una métrica que te da Tienda Nube. ¿Cómo es? ¿Te llega por mail o solamente por Tienda Nube?

\noindent\textbf{Delfina Fernández:} Es la parte de inicio en Tienda Nube. Hay como un cuadrito con cuatro cosas, y podés poner el período: ver ese cuadro basado en hoy, ayer o la semana entera. Ahí es donde dice visitas únicas en el día. Yo siempre lo tengo puesto por el día, pero se puede ver por otros períodos también.

\noindent\textbf{Micaela Palomino:} ¿Y esto te sirve para algo? ¿Después podés llevar esto a una acción?

\noindent\textbf{Delfina Fernández:} Yo creo que sí. Todavía no sé si lo apliqué, pero por ejemplo, me pasó de subir un posteo con un modelo nuevo de bordado, y ese día tuve como diecisiete vistas, que fue un montón. Y después nunca más volví a ver tantas. En cambio, el otro día volví a subir un posteo con las del Mundial, y ahí otra vez subió bastante. Entonces dije: quizás le tengo que meter mucho más a los posteos, y las historias también generan bastante tráfico. Voy tanteando.

\noindent\textbf{Micaela Palomino:} ¿Y después hubo conversión más allá de las visitas? ¿Tuviste más visitas y aumentaron las ventas, o no necesariamente?

\noindent\textbf{Delfina Fernández:} Eso lo estaba pensando el otro día. No sé si va muy de la mano, porque por ejemplo, ese día que tuve diecisiete visitas quizás no tuve venta o tuve una sola. Y el otro día, con el posteo del Mundial, tuve nueve visitas y tuve dos o tres ventas en un día. No sé cómo va eso de la mano.

\noindent\textbf{Tomás Gonçalves:} La cantidad de ventas está atada a lo que haya en el momento, más o menos. Con el Mundial te vino re bien, entonces.

\noindent\textbf{Delfina Fernández:} Sí, re contra. Pero es que el modelo anterior tuve diecisiete visitas y no me compró nadie. En cambio, con este tuve nueve visitas, que fue mucho menos, y tuve tres ventas.

\noindent\textbf{Micaela Palomino:} ¿Qué otras métricas te muestra Tienda Nube además de visitas únicas?

\noindent\textbf{Delfina Fernández:} A ver, me fijo. Te muestra visitas únicas, ventas, facturación y ticket promedio. Se puede filtrar por hoy, ayer y esta semana.

\noindent\textit{[Delfina comparte pantalla mostrando el panel de inicio de Tienda Nube]}

\noindent\textbf{Delfina Fernández:} Este es el inicio. Lo del principio es para mí, lo que tengo que ir marcando: sin ventas por cobrar (porque hay veces que te pagan otro día; Tienda Nube te da hasta tres días para poder pagar por transferencia, y si no, te cancela la compra), cuatro ventas por empaquetar (que son todavía las que no hice), lo mismo con enviar y por retirar. Eso es para ordenarme. Y después están las métricas que les iba diciendo: visitas únicas, ventas, esta semana, ayer. Ustedes díganme qué quieren que aprete, porque yo no tengo idea.

\noindent\textit{--- Seguimiento de ventas y herramientas complementarias ---}

\noindent\textbf{Micaela Palomino:} ¿Vos le das seguimiento a las ventas o solamente lo tenés concentrado acá en Tienda Nube? Por ejemplo, ¿sabés cuál es la remera que más vendiste, o cuál es el diseño que más vendés?

\noindent\textbf{Delfina Fernández:} En Tienda Nube yo sé que si pagás los niveles más pro ---esto lo sé por mi trabajo---, podés ver todo eso. Yo no lo tengo acá porque pago el nivel uno. Igual, yo hago el seguimiento. Antes lo hacía por Tienda Nube: me iba anotando todo. Es más, las ventas que me hacían por boca a boca yo igual las simulaba en Tienda Nube. Pero ahora Facundo me hizo una app, entonces guardo todo ahí; voy marcando todo ahí.

\noindent\textbf{Micaela Palomino:} ¿Y la app que te hizo Facundo, qué seguimiento le das? ¿Para qué métricas la usás?

\noindent\textbf{Delfina Fernández:} Las ventas las cargo todas ahí. Después también me hizo como parte del proceso de producción: yo voy marcando si ya está bordada, si se va a lavar, después a planchar, después a empaquetar y después ya está entregada. Y todo el tema de los ingresos y egresos también los marco ahí. Todas las remeras que compro, las veces que compro, las marco ahí, y después se va haciendo toda la cuenta sola.

\noindent\textit{[Facundo, colaborador de Delfina, muestra brevemente la sección de analytics de la app]}

\noindent\textbf{Facundo (colaborador):} En Analytics tenemos todo lo de los ingresos totales, egresos totales, balance neto por mes y por semana. Después también está la parte de ver cuál es la que más se vende.

\subsubsection*{Bloque 2 · Gestión de inventario y stock}

\noindent\textit{--- Publicidad y temor al quiebre de stock ---}

\noindent\textbf{Tomás Gonçalves:} ¿Hacés publicidad?

\noindent\textbf{Delfina Fernández:} Todavía no. Ese es mi próximo paso. Literalmente esta semana estoy tratando de organizarme para ver qué necesito hacer primero. El tema es que yo tardo en hacer cada producto. Ahora, ponele, ya sé cuál es mi límite de ventas que puedo tener por semana: son las remeras que yo llego a hacer. Si no, las entrego tarde. Ahora estoy en el límite, por suerte. Entonces, si pongo publicidad, un poco de miedo me da de que se vaya todo a la mierda.

\noindent\textbf{Tomás Gonçalves:} Alto miedo. ¿Cómo es el tema de las publicidades? Me imagino que tenés que estar bien preparada.

\noindent\textbf{Delfina Fernández:} Claro, porque en el otro lugar donde trabajo, antes pasaba eso: era activar publicidad y a los dos días darla de baja porque decíamos que ya no dábamos abasto. Entonces un poco de miedo tengo.

\noindent\textbf{Tomás Gonçalves:} Es como que tenés que estar pendiente del stock y todo eso.

\noindent\textbf{Delfina Fernández:} Sí, eso es lo que me mata, porque todavía no tengo un stock tremendo. Entonces me tengo que estar muy encima.

\noindent\textit{--- Quiebre de stock y reposición de insumos ---}

\noindent\textbf{Micaela Palomino:} ¿Con los modelos de bordado te pasa lo mismo de guiarte por intuición?

\noindent\textbf{Delfina Fernández:} Los modelos de bordado no; eso sí me doy cuenta más por lo que pidieron. Por ejemplo, yo sé que el pimiento y las aceitunitas se venden mucho. Llegó un momento en que me estaba quedando sin mostacillas rojas para el pimiento y dije: ¿qué hago ahora? No las conseguía, no las conseguía. Me compré una bolsa grande, y ahora nadie me pide más el pimiento porque vino el invierno. Pero sí, eso también lo veo por la app de Facundo, que ahí tengo el orden de las que más fui vendiendo. Por ejemplo, ahora me estoy quedando sin mostacillas de otro color que no vendo tanto, y digo: bueno, ya está.

\noindent\textbf{Micaela Palomino:} Si hubieses sabido que la de pimiento se iba a vender más, ¿hubieses comprado de antemano más mostacillas? ¿Te pasó alguna vez que vendiste más de lo que pensabas?

\noindent\textbf{Delfina Fernández:} Sí, me pasó justamente con el pimiento y con las aceitunitas. Me quedé sin mostacillas y dije: tengo que salir ya a comprar, porque no tenía, y no pensé que se iban a vender tanto. De la nada, en una o dos semanas, las aceitunitas era la única que me pedían.

\noindent\textbf{Tomás Gonçalves:} Claro, te enterás en el momento cuando ya se agotó.

\noindent\textbf{Delfina Fernández:} Claro, yo voy tanteando cuánto tengo, pero es como que nunca se sabe qué venta va a entrar.

\noindent\textbf{Tomás Gonçalves:} Entonces, proyectar cuántas ventas vas a tener la semana que viene es imposible.

\noindent\textbf{Delfina Fernández:} Con los modelos anteriores, que tenía varios y eran todos de verano, sí, era más difícil. Ahora que estoy con las del Mundial, hice tres modelos. Hay uno que me lo pidieron una sola vez; los otros dos están muy peleados. Igual, sé que me pedían mucho más la de Argentina que la de las Carapela, por ejemplo. Pero no te puedo decir con métricas; es más intuición. También me doy cuenta cuando subo historias o posteos a Instagram y me contestan bastante: me contestan más una que la otra, y voy por ahí.

\subsubsection*{Bloque 3 · Análisis de ventas y decisiones comerciales}

\noindent\textit{--- Patrones de venta y estacionalidad ---}

\noindent\textbf{Micaela Palomino:} ¿Sabés qué día es el que más vendés?

\noindent\textbf{Delfina Fernández:} No, eso por ahora está como muy difuso. Yo creo que ni la app lo podría captar, porque es muy variable. Por ejemplo, vendo más del boca en boca: voy a Pilates y llevo una para que alguien se pruebe, y ya ahí vendo tres. Entonces quizás te puedo decir que los jueves que voy a Pilates vendo, pero de repente un jueves subo un posteo y la gente me compra tres en un día también por Tienda Nube. Por ahora no sé si hay tantas ventas como para decir que un día es el que más compran. Sí puedo decir horarios: me compran más a la mañana, entre las ocho y las once, o a la noche, tipo ocho o nueve. Es como que la gente o termina el día y se pone a comprar, o se levanta y se pone a comprar.

\noindent\textbf{Tomás Gonçalves:} Con el Mundial, ¿vos dijiste ``voy a lanzar un producto relacionado a eso'', o sacás productos random?

\noindent\textbf{Delfina Fernández:} Los modelos nuevos los venía sacando random: si encontraba alguno que me gustaba y me emocionaba hacerlo, lo hacía y lo sacaba. Pero con las del Mundial dije: sí, algo tengo que hacer, porque se va a vender. Este año igual hay cero emoción de Mundial, como se habrán dado cuenta, pero vienen bien. Tengo tres modelos, varían los tipos de remera también, y se viene vendiendo bien. Fue tipo una cápsula mundialista.

\noindent\textit{--- Política de precios ---}

\noindent\textbf{Micaela Palomino:} ¿Cómo van los precios? ¿Te basás en algo? ¿Los ajustás?

\noindent\textbf{Delfina Fernández:} Eso es un tema que me cuesta mucho. Para arrancar, en noviembre, sacamos los costos (lo hizo Facundo). Yo no quería que fuera una prenda cara, porque sé que siendo un Instagram de trescientos seguidores, y aparte al principio siempre te compran los conocidos, me cuesta ponerle un precio alto. Los precios los mantengo lo más que puedo por ahora; los puse desde el principio y no se movieron mucho. Cuando fui sumando otro tipo de remeras que están más caras, ahí sí o sí esa diferencia se la sumo, porque si no me la como yo. No es lo mismo una musculosa que una remera boxy fit.

\noindent\textit{--- Datos concretos versus intuición ---}

\noindent\textbf{Micaela Palomino:} Hoy en día, ¿cuántas de las decisiones que tomás para el emprendimiento están basadas en datos concretos y cuántas en la intuición y experiencia? Por ejemplo, cuando tenés que hacer una compra de remeras, ¿comprás porque tenés datos de que se van a vender, o más por intuición?

\noindent\textbf{Delfina Fernández:} Un poco y un poco. Por ejemplo, ya sé que el talle M es el que más me piden, entonces ahí sí sé por qué: es el que más pidieron. Y a la vez, si alguien quería un S, le ofrezco el M porque sé que no es una gran diferencia, y lo mismo con el L. Después, el negro no se vendió tanto como la blanca, porque también tengo muchos más modelos de bordado que van más con la blanca que con la negra. Pero sí, es un poco y un poco por ahora.

\subsubsection*{Bloque 4 · Puntos de dolor y visibilidad del negocio}

\noindent\textbf{Tomás Gonçalves:} A mí me asustaría la incertidumbre: no saber qué vendo, qué no vendo.

\noindent\textbf{Delfina Fernández:} Sí, también te cuento: hice una promoción con las del Mundial. Dije que si me compran esta remera y ganamos la cuarta, les guardo la cuarta estrella gratis. Pero solo una persona me la compró. Tampoco sé si me conviene, porque es la que más tardo en hacer. Así que lo voy a pensar.

\subsubsection*{Bloque 5 · Validación de la propuesta de valor}

\noindent\textit{--- Recomendaciones de reposición ---}

\noindent\textbf{Tomás Gonçalves:} Imaginá que Tienda Nube te ofrece una opción que, en base a tu historial, te dice: ``Te conviene reponer X producto antes del fin de semana porque el sábado vas a tener unas ventas esperadas de veinte unidades.'' ¿Te coparía tener algo así?

\noindent\textbf{Delfina Fernández:} Eso estaría buenísimo, porque más que nada a mí me sirve por el material. Nunca sé si comprar más o si me los voy a quedar sin vender. Entonces me sirve por eso, porque si termino comprando algo que no lo voy a vender, es pérdida. Eso sería buenísimo.

\noindent\textit{--- Visibilidad por producto ---}

\noindent\textbf{Micaela Palomino:} Vos lo que decías era que te guiabas por la intuición, pero en realidad la intuición es la interacción que tenés en Instagram. Después nos contabas que en Tienda Nube te aparecen las visitas únicas, pero no te dice específicamente de qué producto. Si te dijera que hoy hiciste un posteo y tuviste quince visualizaciones, y las quince fueron de determinado producto, tal vez es un índice de que ese producto probablemente lo vendas.

\noindent\textbf{Delfina Fernández:} Sí, eso estaría buenísimo también. Estaría muy bueno.

\noindent\textbf{Delfina Fernández:} Ahora estoy con el primer nivel pago de Tienda Nube. Estaba con el que es gratis y no te daba nada de métricas. Recién el mes pasado lo empecé a pagar.

\noindent\textit{--- Descripción de la propuesta y reacciones ---}

\noindent\textbf{Tomás Gonçalves:} La herramienta que más o menos te describimos, ¿qué valor le darías? ¿Cuánto pagarías por algo así?

\noindent\textbf{Delfina Fernández:} ¿Solo por eso, o si viene incluido en Tienda Nube?

\noindent\textbf{Micaela Palomino:} Lo que pensábamos era armar un agente que se conecte a Tienda Nube, pero que ofrezca acciones concretas. No eso de ``tuviste cinco visitas'' y vos decís: ¿qué hago ahora? Sino más bien: aumentá el stock de esto, dejá de comprar mostacillas de las tres estrellas porque nadie entró a verlas, los días jueves son los que más vendés (o al revés: los miércoles son los que menos vendés, entonces poné un descuento para que los miércoles la gente entre). Datos más concretos, acciones en vez de métricas. Decirte exactamente lo que tenés que hacer para aumentar lo que vos quieras, dependiendo de tu objetivo: puede ser aumentar las ventas, o quizás sos más conservadora y preferís vender todo el stock. Pensamos dos formas: una era que se conecte directamente a Tienda Nube, porque entendemos que la gente quizás no quiere usar dos herramientas.

\noindent\textbf{Delfina Fernández:} No, y aparte es que los clientes entran a la página online, que es Tienda Nube. Entonces está bueno que agarre toda la información de ahí, porque ahí es donde se vende.

\noindent\textbf{Micaela Palomino:} La otra opción era armar una plataforma separada que haga lo mismo. Lo estamos pensando.

\noindent\textbf{Delfina Fernández:} Está buenísimo. Con todos estos datos hay infinidad de cosas que podés hacer: desde poner descuentos, meter publicidad solo con esa remera, subir un millón de posteos ese día. Está buenísimo. Y aparte, sirve, porque yo quizás veo las diecisiete vistas y digo: ay, qué bueno, se re pegó este posteo. Después hubo otro igual o parecido y no se pega, y no tengo esas vistas. Entonces digo: ¿qué hago?

\noindent\textbf{Micaela Palomino:} Siempre es el mismo problema: no entendés por qué fueron esas quince vistas. Si te dijéramos que fueron por tal cosa, ya sabés qué patrón repetir o qué posteo funciona.

\noindent\textbf{Delfina Fernández:} Sí, y también de dónde vienen las visitas. Quizás vienen del perfil, o quizás vienen de la historia que subí. Eso está buenísimo.

\noindent\textbf{Micaela Palomino:} Por ejemplo, si alguien comparte tu link por WhatsApp, no es lo mismo que entre por un posteo tuyo. Sabés que funciona el boca en boca, y después sabés que tenés más conversión por los posteos, entonces decís: le tengo que dedicar un poquito más a los posteos.

\noindent\textit{--- Disposición a pagar ---}

\noindent\textbf{Delfina Fernández:} Bueno, contestando a la pregunta de cuánto pagaría: yo ahora estoy pagando veintisiete mil pesos. Yo creo que entre cuarenta y cincuenta mil pagaría si viene con esto, porque es lo que me va a terminar de dar la devolución de lo que estás haciendo; es lo que me estaría faltando ahora. Así que sí, yo creo que lo vale. No me acuerdo cuánto está el próximo plan de Tienda Nube; no sé si estaba como setenta u ochenta mil. No sé qué tiene tampoco, no me puse a fijar. Pero creo que es un muy buen extra.

\subsubsection*{Bloque 6 · Cierre}

\noindent\textbf{Micaela Palomino:} Más allá de todo lo que te contamos que queremos hacer, si te damos la posibilidad de decirnos qué te gustaría que hagamos: ¿qué otras funcionalidades te gustaría que tenga?

\noindent\textbf{Delfina Fernández:} Es muy difícil para contestar ahora.

\noindent\textbf{Micaela Palomino:} ¿Qué es lo que más te cuesta hoy, que decís: por favor, resuélvanme esto?

\noindent\textbf{Delfina Fernández:} Quizás el tema publicidad. Todo lo que me venís nombrando es literalmente lo que me venía haciendo ruido, porque no sé cuándo, es todo incertidumbre: no sé cuánto voy a vender, cómo, qué. Quizás que me ayude con el tema publicidad: ``tengo floja esta remera'', que me lo avise y yo le pongo toda la publicidad del mundo a esa remera. O que me prediga cuánto voy a vender de tal modelo esta semana, que puede ser como el boom; así subo la publicidad al máximo o la bajo. Eso está buenísimo. Ahora no se me ocurre algo extra, porque siento que de verdad está muy bueno lo que me estaban contando. No sé si le falta algo por ahora.

\noindent\textbf{Micaela Palomino:} Siento que sería muy bueno que, en vez de recibir el pedido para ponerte a hacer el diseño y el bordado, lo podrías hacer de antemano y ya los tendrías preparados para cuando te lo pidan.

\noindent\textbf{Delfina Fernández:} Sí, eso sería buenísimo. Aparte hay baches de ventas. Antes del Mundial estaba medio sin saber qué hacer y me puse a hacer un par de antemano. Ahora me resolvió la vida, porque de la nada en una semana vendí diez o doce, y digo: las tengo que hacer todas para la semana que viene, pero ya tenía un par hechas. Eso resuelve de muchas maneras.

\noindent\textbf{Tomás Gonçalves:} Por ahora está muy completo. Pero capaz que nosotros nos salteamos algo. ¿Hay algo que quieras saber o que te gustaría que hayamos preguntado?

\noindent\textbf{Delfina Fernández:} No tenía expectativas; es más, tenía un poco de miedo. Si me preguntan mucho de métricas y esas cosas que yo todavía no me acostumbro a usar, porque ni siquiera pago publicidad, dije: ay, ¿qué voy a hacer? Pero creo que estuvo bastante bien. Me gustó mucho lo que me están contando de qué van a hacer. El concepto ya me encantó: esto de que se va a acabar la incertidumbre, que es lo que te mata un poco.

\noindent\textbf{Delfina Fernández:} Se me ocurre también poder linkearlo con Instagram. Quizás como poder ver: este posteo está hecho de tal y tal manera, y esto es lo que funcionó. Sacar datos de ahí también, para decir: esto vendió.

\noindent\textbf{Delfina Fernández:} Yo ahora me estoy manejando mucho más con Instagram que con Tienda Nube, porque no termino de usar Tienda Nube bien. Dice cosas y digo: ¿qué hago con esto? En Instagram es todo un tema. Todavía no pago publicidad, y estaría bueno que si pago esto de Tienda Nube, tenerlo linkeado con Instagram y ver qué sale de ahí.

\noindent\textbf{Micaela Palomino:} Muchas gracias por el aporte y por conectarte un domingo a la noche.

\noindent\textbf{Delfina Fernández:} Cualquier cosa me dicen, yo no tengo problema.

\noindent\textbf{Tomás Gonçalves:} Muchas gracias por todo.

\begin{center}\textit{--- Fin de la entrevista ---}\end{center}
